\begin{abstract}
Approximate nearest-neighbour indexes are usually compared at a recall target, but applications see the task quality of
the final ranking. We compare a binary cascade (512-bit sign codes, an HNSW walk or exact scan, an int8 screen, float16
rescoring) with indexes from FAISS, USearch, hnswlib, ScaNN, DiskANN and Intel SVS, each at the fastest configuration
that keeps 99\% of the task quality of exact search, chosen on dev queries and frozen before any held-out run. In a
pre-registered replication on two untouched sets (CQADupStack, 457k documents; MS MARCO, 1M passages) on three cloud
CPUs, the cascade was faster than the rival chosen on dev in all \nm{rep-n} tests (\nm{rep-min}--\nm{rep-max} times,
Holm-corrected); one pass depended on a page-fault rule that in practice dropped only our jobs, and a post hoc rerun with
the index preloaded passes all six with no block dropped. Quality non-inferiority held on MS MARCO only. An exploratory
main study (seven code and BEIR sets, 7k to 1M vectors, three models, five CPUs) adds two findings. On the same codes at
equal graph settings, our walk is \nm{post-eqb-min}--\nm{post-eqb-max} times faster than the binary HNSW of FAISS and
USearch while computing about as many distances; hardware counters place the difference in instructions and branch
misses. On the code sets, points with recall@10 as low as \nm{recdoc-min} keep 99\% of exact quality; on the replication
sets recall stays above 0.9 at that quality but still misjudges it. The cascade loses below about 20k vectors, uses more
memory, builds slowly and, with its default memory-mapped index, depends on the page cache. We release the harness, the
frozen configurations and every record.
\end{abstract}
